\chapter{Zinc Deficiency}
\index{Zinc Deficiency}
\label{sec:Zinc Deficiency}

\section{Overview}
Zinc deficiency is a mineral deficiency disorder common in low-resource settings due to low dietary intake of animal-source foods and high phytate content (inhibits zinc absorption). At-risk groups include pregnant and lactating women, children (especially with diarrhea), the elderly, vegetarians, and those with malabsorption (acrodermatitis enteropathica---autosomal recessive zinc transporter defect). This metabolic disorder involves dysregulation of normal bodily processes, often requiring ongoing monitoring and management. It is increasingly common due to lifestyle and dietary factors. Metabolic disorders involve disruption of normal biochemical processes essential for health. These conditions may result from enzyme deficiencies, hormonal imbalances, or lifestyle factors. Many metabolic disorders are chronic and require ongoing management. Lifestyle modifications are often the cornerstone of treatment.
\section{Causes}
This condition happens because of deficiency of zinc, an essential trace element that serves as a cofactor for over 300 enzymes involved in growth, immune function, protein synthesis, wound healing, and DNA synthesis. Metabolic disorders arise from dysregulation of normal biochemical pathways. This may result from enzyme deficiencies, hormonal imbalances, or lifestyle factors that overwhelm normal homeostatic mechanisms.   
\section{Risk Factors}
\begin{itemize}[noitemsep]
  \item Genetic predisposition and family history
  \item Advancing age
  \item Lifestyle factors (diet, exercise, smoking, alcohol)
  \item Environmental exposures
  \item Underlying medical conditions
  \item Medication use
\end{itemize}
\section{Signs and Symptoms}
\subsection{Common symptoms}
\begin{itemize}[noitemsep]
  \item You may experience growth retardation, delayed sexual maturation, hypogonadism, impaired immune function (recurrent infections, especially diarrheal diseases), alopecia, dermatitis (periorificial and acral, vesiculobullous or eczematous---characteristic of acrodermatitis enteropathica), poor wound healing, night blindness, anorexia, dysgeusia, and neuropsychiatric changes.
  \item Pain or discomfort in the affected area
  \end{itemize}
\subsection{Emergency warning signs}
\begin{itemize}[noitemsep]
  \item Severe or worsening symptoms of zinc deficiency require immediate medical evaluation
\end{itemize}
\section{Progression and Stages}
The condition may progress through different stages of severity over time. Early stages often involve mild or subtle symptoms that may be overlooked. Moderate stages involve more noticeable symptoms affecting daily function. Advanced stages may involve significant impairment and complications.
\section{Complications}
\begin{itemize}[noitemsep]
  \item Disease progression leading to worsening symptoms
  \item Secondary infections or injuries
  \item Side effects of treatment
  \item Reduced quality of life
  \item Emotional and psychological impact
\end{itemize}
\section{Diagnosis}
Doctors diagnose this using low plasma zinc (less than 60--70 \textmu g/dL).

\begin{itemize}[noitemsep]
  \item Comprehensive medical history and physical examination
  \item Laboratory tests to assess organ function
  \item Imaging studies to visualize affected structures
  \item Specialized testing depending on the affected system
  \item Consultation with relevant specialists
\end{itemize}
\section{Prevention}
\begin{itemize}[noitemsep]
  \item Maintain a healthy lifestyle with balanced diet and exercise
  \item Avoid known risk factors
  \item Attend regular medical check-ups for early detection
  \item Follow recommended screening guidelines
\end{itemize}
\section{Treatment Options}
\begin{itemize}[noitemsep]
  \item Treatment focuses on zinc sulfate or acetate 1--3 mg/kg/day (adults: 25--50 mg elemental zinc daily)
  \item For acrodermatitis enteropathica, lifelong supplementation is required.
  \item Medications to manage symptoms and address underlying mechanisms
  \item Lifestyle modifications and self-care measures
  \item Physical therapy and rehabilitation
  \item Surgical intervention when indicated
  \item Regular monitoring and follow-up care
\end{itemize}
\section{Living with It}
\begin{itemize}[noitemsep]
  \item Follow your treatment plan as prescribed
  \item Maintain regular follow-up with your healthcare provider
  \item Make necessary lifestyle adjustments
  \item Seek support from family, friends, or support groups
  \item Stay informed about your condition
\end{itemize}
\section{Outlook}
Most people recover well with supplementation.
